\chapter*{Abstract}

This report develops data-driven indicators for the predictive maintenance of a metro
train's pneumatic system from the MetroAT dataset: one year of 1\,Hz, min-max-normalised
sensor recordings from a single Wiener Linien train. The analysis is organised around one
methodological principle: the avoidance of leakage, or circularity. Operational states
are defined from motion alone, so that the air-system sensors, which never entered that
definition, remain independent evidence.

Clustering the motion features recovers four operational states: standing (48\%),
deceleration (20\%), cruising (17\%), and accelerating (14\%). The cluster count $k=4$ was
selected by the largest information-criterion improvement and confirmed stable under
resampling. Restricting the state definition to motion corrected an earlier error in which
a stationary train holding its brake was mislabelled as deceleration; the corrected
deceleration state is 99.9\% genuine. On held-out data, the air-system sensors recover the
deceleration state with a ROC-AUC of 0.95, and 0.92 even when the brake command is withheld
and only independent health sensors are used, so deceleration leaves a fingerprint across
the pneumatic system.

Deceleration intensity was tested for class structure and found to be a continuum, not a
set of discrete classes. A gradient-boosted tree explains a moderate share of intensity
from the non-speed sensors out-of-sample ($R^2 \approx 0.43$--$0.47$, against a
guess-the-average baseline), and about 83\% of the specific energy dissipated per stop.
Finally, braking behaviour was tested for a pre-failure signal ahead of documented
brake-system failures. Across effect-size, change-point, and energy analyses the result is
near-null; with only eight brake failures the study is exploratory and provides no
deployable failure detector. The negative result suggests that brake faults may manifest
during charging and idle periods rather than during deceleration, which is the direction a
follow-up should take.
